% GENERATED FILE. Edit canonical lessons, then run npm run generate:books.
% canonical-source-hash: fnv1a64:c5739c3beac8f668
% canonical-lessons: SA-C405-javah, SA-C405-satvaram, SA-C405-pradhavati, SA-C405-antimah, SA-C405-ayati

\chapter{Hurry and Speed}
\label{ch:sa-hurry-and-speed}
\hlchaptermodality{405}

\begin{chapteropening}
\textbf{By the end of this chapter:} \emph{I can say who is in a hurry and who arrives last.}

Complete the last lesson of chapter 405: I can say who is in a hurry and who arrives last.
\end{chapteropening}

\section[\texorpdfstring{\sk{जवः} (javaḥ)}{जवः (javaḥ)}]{\sk{जवः} (javaḥ) — swiftness}

\label{lesson:SA-C405-javah}



\begin{quote}
Before the new one: say the Sanskrit for takes cover, then the Sanskrit for deserted.
\end{quote}

\subsection*{You'll want to know: \sk{जवः}}
\textbf{\sk{जवः}} — \emph{javaḥ} — \textquotedblleft{}swiftness\textquotedblright{}.

\begin{quote}
\sk{अश्वस्य} \sk{जवः} \sk{महान्} \sk{अस्ति।}
\end{quote}

\begin{quote}
\emph{aśvasya javaḥ mahān asti.}
\end{quote}

\begin{quote}
The horse's swiftness is great.
\end{quote}

A masculine noun.

\begin{cousinweb}
\textbf{\sk{जवः}} (\emph{javaḥ}) comes from the root √ju, \textquotedblleft{}to hasten\textquotedblright{}; MW gives \emph{java} as speed, velocity, swiftness (MW).
\end{cousinweb}

\begin{grammarlens}[title={What you've built}]
One more word for getting around.
\end{grammarlens}

\subsection*{Your turn}
\begin{itemize}
\raggedright
  \item \emph{Say it:} \emph{javaḥ}
  \item \emph{Say it:} \emph{javaḥ}, once more
  \item \emph{Say it:} \emph{aśvasya javaḥ mahān asti.} — The horse's swiftness is great.
  \item \emph{Recall:} say the Sanskrit for takes cover, then the Sanskrit for deserted, then say \emph{javaḥ} again
\end{itemize}

\begin{tcolorbox}[breakable,colback=teal!4,colframe=teal!35!black,title={Before you move on}]
What does \sk{जवः} mean? (\textquotedblleft{}Swiftness\textquotedblright{}.) Say it once more.
\end{tcolorbox}
\section[\texorpdfstring{\sk{सत्वरम्} (satvaram)}{सत्वरम् (satvaram)}]{\sk{सत्वरम्} (satvaram) — in a hurry; hastily}

\label{lesson:SA-C405-satvaram}



\begin{quote}
Before the new one: say the Sanskrit for deserted, then the Sanskrit for swiftness.
\end{quote}

\subsection*{You'll want to know: \sk{सत्वरम्}}
\textbf{\sk{सत्वरम्}} — \emph{satvaram} — \textquotedblleft{}in a hurry; hastily\textquotedblright{}.

\begin{quote}
\sk{यात्री} \sk{सत्वरं} \sk{स्थानकं} \sk{गच्छति।}
\end{quote}

\begin{quote}
\emph{yātrī satvaraṃ sthānakaṃ gacchati.}
\end{quote}

\begin{quote}
The passenger goes to the station in a hurry.
\end{quote}

\begin{cousinweb}
\textbf{\sk{सत्वरम्}} (\emph{satvaram}) is built from \emph{sa}, \textquotedblleft{}with\textquotedblright{}, and \emph{tvarā}, \textquotedblleft{}haste\textquotedblright{}, from the root √tvar, \textquotedblleft{}to hurry\textquotedblright{}; MW gives \emph{satvara} as making haste and its neuter as an adverb (MW).
\end{cousinweb}

\begin{grammarlens}[title={What you've built}]
One more word for getting around.
\end{grammarlens}

\subsection*{Your turn}
\begin{itemize}
\raggedright
  \item \emph{Say it:} \emph{satvaram}
  \item \emph{Say it:} \emph{satvaram}, once more
  \item \emph{Say it:} \emph{yātrī satvaraṃ sthānakaṃ gacchati.} — The passenger goes to the station in a hurry.
  \item \emph{Recall:} say the Sanskrit for deserted, then the Sanskrit for swiftness, then say \emph{satvaram} again
\end{itemize}

\begin{tcolorbox}[breakable,colback=teal!4,colframe=teal!35!black,title={Before you move on}]
What does \sk{सत्वरम्} mean? (\textquotedblleft{}In a hurry; hastily\textquotedblright{}.) Say it once more.
\end{tcolorbox}
\section[\texorpdfstring{\sk{प्रधावति} (pradhāvati)}{प्रधावति (pradhāvati)}]{\sk{प्रधावति} (pradhāvati) — dashes off; rushes forward}

\label{lesson:SA-C405-pradhavati}



\begin{quote}
Before the new one: say the Sanskrit for swiftness, then the Sanskrit for in a hurry.
\end{quote}

\subsection*{You'll want to know: \sk{प्रधावति}}
\textbf{\sk{प्रधावति}} — \emph{pradhāvati} — \textquotedblleft{}dashes off; rushes forward\textquotedblright{}.

\begin{quote}
\sk{बालः} \sk{लोकयानं} \sk{प्रति} \sk{प्रधावति।}
\end{quote}

\begin{quote}
\emph{bālaḥ lokayānaṃ prati pradhāvati.}
\end{quote}

\begin{quote}
The boy dashes off towards the bus.
\end{quote}

Past (he …): \sk{सः} \sk{लोकयानं} \sk{प्रति} \sk{प्रधावितवान्}, \emph{saḥ lokayānaṃ prati pradhāvitavān}.

\begin{cousinweb}
\textbf{\sk{प्रधावति}} (\emph{pradhāvati}) is built from \emph{pra}, \textquotedblleft{}forward\textquotedblright{}, and the root √dhāv, \textquotedblleft{}to run\textquotedblright{}, which also gives \emph{dhāvati}, \textquotedblleft{}runs\textquotedblright{}; MW gives \emph{pra-dhāv} as to run forward, run off or rush, with the participle \emph{pradhāvita}, \textquotedblleft{}run off, set out\textquotedblright{} (MW).
\end{cousinweb}

\begin{grammarlens}[title={What you've built}]
One more word for getting around.
\end{grammarlens}

\subsection*{Your turn}
\begin{itemize}
\raggedright
  \item \emph{Say it:} \emph{pradhāvati}
  \item \emph{Say it:} \emph{pradhāvati}, once more
  \item \emph{Say it:} \emph{bālaḥ lokayānaṃ prati pradhāvati.} — The boy dashes off towards the bus.
  \item \emph{Recall:} say the Sanskrit for swiftness, then the Sanskrit for in a hurry, then say \emph{pradhāvati} again
\end{itemize}

\begin{tcolorbox}[breakable,colback=teal!4,colframe=teal!35!black,title={Before you move on}]
What does \sk{प्रधावति} mean? (\textquotedblleft{}Dashes off; rushes forward\textquotedblright{}.) Say it once more.
\end{tcolorbox}
\section[\texorpdfstring{\sk{अन्तिमः} (antimaḥ)}{अन्तिमः (antimaḥ)}]{\sk{अन्तिमः} (antimaḥ) — last; final}

\label{lesson:SA-C405-antimah}



\begin{quote}
Before the new one: say the Sanskrit for in a hurry, then the Sanskrit for dashes off.
\end{quote}

\subsection*{You'll want to know: \sk{अन्तिमः}}
\textbf{\sk{अन्तिमः}} — \emph{antimaḥ} — \textquotedblleft{}last; final\textquotedblright{}.

\begin{quote}
\sk{अन्तिमं} \sk{लोकयानं} \sk{रात्रौ} \sk{गच्छति।}
\end{quote}

\begin{quote}
\emph{antimaṃ lokayānaṃ rātrau gacchati.}
\end{quote}

\begin{quote}
The last bus leaves at night.
\end{quote}

\begin{cousinweb}
\textbf{\sk{अन्तिमः}} (\emph{antimaḥ}) is built on \emph{anta}, \textquotedblleft{}an end\textquotedblright{}, with the ending \emph{-ima}; MW gives \emph{antima} as final, ultimate, last (MW).
\end{cousinweb}

\begin{grammarlens}[title={What you've built}]
One more word for getting around.
\end{grammarlens}

\subsection*{Your turn}
\begin{itemize}
\raggedright
  \item \emph{Say it:} \emph{antimaḥ}
  \item \emph{Say it:} \emph{antimaḥ}, once more
  \item \emph{Say it:} \emph{antimaṃ lokayānaṃ rātrau gacchati.} — The last bus leaves at night.
  \item \emph{Recall:} say the Sanskrit for in a hurry, then the Sanskrit for dashes off, then say \emph{antimaḥ} again
\end{itemize}

\begin{tcolorbox}[breakable,colback=teal!4,colframe=teal!35!black,title={Before you move on}]
What does \sk{अन्तिमः} mean? (\textquotedblleft{}Last; final\textquotedblright{}.) Say it once more.
\end{tcolorbox}
\section[\texorpdfstring{\sk{आयाति} (āyāti)}{आयाति (āyāti)}]{\sk{आयाति} (āyāti) — arrives; comes}

\label{lesson:SA-C405-ayati}



\begin{quote}
Before the new one: say the Sanskrit for dashes off, then the Sanskrit for last.
\end{quote}

\subsection*{You'll want to know: \sk{आयाति}}
\textbf{\sk{आयाति}} — \emph{āyāti} — \textquotedblleft{}arrives; comes\textquotedblright{}.

\begin{quote}
\sk{अन्तिमं} \sk{लोकयानं} \sk{विलम्बेन} \sk{आयाति।}
\end{quote}

\begin{quote}
\emph{antimaṃ lokayānaṃ vilambena āyāti.}
\end{quote}

\begin{quote}
The last bus arrives late.
\end{quote}

Past (it …): \sk{लोकयानम्} \sk{आयातवत्}, \emph{lokayānam āyātavat}.

\begin{cousinweb}
\textbf{\sk{आयाति}} (\emph{āyāti}) is built from \emph{ā}, \textquotedblleft{}towards\textquotedblright{}, and the root √yā, \textquotedblleft{}to go\textquotedblright{}; MW gives \emph{ā-yā} as to come near, to arrive, with the participle \emph{āyāta} (MW).
\end{cousinweb}

\begin{grammarlens}[title={What you've built}]
One more word for getting around.
\end{grammarlens}

\subsection*{Your turn}
\begin{itemize}
\raggedright
  \item \emph{Say it:} \emph{āyāti}
  \item \emph{Say it:} \emph{āyāti}, once more
  \item \emph{Say it:} \emph{antimaṃ lokayānaṃ vilambena āyāti.} — The last bus arrives late.
  \item \emph{Recall:} say the Sanskrit for dashes off, then the Sanskrit for last, then say \emph{āyāti} again
\end{itemize}

\begin{tcolorbox}[breakable,colback=teal!4,colframe=teal!35!black,title={Before you move on}]
What does \sk{आयाति} mean? (\textquotedblleft{}Arrives; comes\textquotedblright{}.) Say it once more.
\end{tcolorbox}
